\bookchapter{第二十三章：プログラミング的思考とオートメーション入門}

\chapterintro{前二章では、スマートフォンに安全な「境界壁」を築き、壁の中で日に日に賢くなる「AIの頭脳」を見てきました。ここで、スマートフォンの究極のポテンシャルの一つを解放しましょう：プログラマーのように「考え」、設定した連続タスクを自動実行させること。「プログラミングなんてできないし、自分には関係ない」と思うかもしれません。安心してください、本章では一行のコードも書く必要はありません——現代のスマートフォンは「プログラミング」のハードルを、ドラッグ＆ドロップで済むレベルまで下げています。}

\booksection{一、オートメーションとは何か？}

例えば：毎朝起きた時、手動で「おやすみモード」をオフにし、WeChatを開いてメッセージを確認し、天気アプリを開いて今日の服装を決め、NetEase Cloud Musicを開いて曲を再生する。これらのステップを一つのオートメーションフローにまとめ、「毎朝7時に自動実行」と設定すれば、起きてスマートフォンを手に取るだけで、すべてが準備完了です。想像するだけで素晴らしいと思いませんか？\textbf{オートメーションの本質は「もし……なら……」という論理チェーン、プログラミングにおけるif-elseです。この考え方を理解すれば、一生コードを書かなくても、問題処理の条理が格段に良くなります。核心となる要素：}トリガー条件

：いつ開始するか？「毎朝7時」「自宅のWi-Fiに接続した時」「特定の場所に到着した時」「バッテリー残量が20\%を下回った時」などが設定できます。

アクション

\begin{enumerate}
  \item \textbf{：具体的に何をするか？WeChatメッセージの送信、特定のアプリを開く、音量調整、壁紙の切り替え、音楽再生、通知の送信……スマートフォンで手動できることのほぼすべてを、オートメーションで実行できます。}中間論理
  \item \textbf{（任意だが強力）：状況AならXを実行、そうでなければYを実行。例えば「平日ならニュースポッドキャストを再生、週末なら音楽を再生」など、これを条件分岐と呼びます。}二、iOS「ショートカット」——Appleユーザーのオートメーション最強ツール
  \item \textbf{Appleの「ショートカット」（Shortcuts）アプリは標準搭載されており、ダウンロード不要です。その核心概念は、各種操作を積み木のように組み合わせることです。まずいくつかの重要な概念を説明します：}「ショートカット」
\end{enumerate}

\booksection{二、iOS「ショートカット」——Appleユーザーのオートメーション最強ツール}

「オートメーション」

\begin{itemize}
  \item \textbf{：トリガー条件により自動実行される指示。手動でトリガーする必要がありません。例えば「スマートフォンがカーブルートゥースに接続した時、自動的に地図アプリを開く」。オートメーションは「オートメーション」タブで設定し、「ショートカット」とは分かれています。}「アクション」
  \item \textbf{：最小の実行単位。アプリを開く、メッセージを送信する、天気を取得する、明るさを調整する——それぞれが一つのアクションです。何百、何千ものアクションを繋げることができます。}実用的な例をいくつか紹介します。実際にやってみれば理解できるはずです：
  \item \textbf{実例1：ワンタップ「おはようフロー」}毎朝タップする（またはSiriに「おはよう」と言う）だけで、スマートフォンが自動的に：今日の天気を取得天気と気温を朗読今日のカレンダー予定を表示よく使う音楽アプリを開いて再生を開始。操作手順：新規ショートカットアクション「現在の天気を取得」を追加「テキストを朗読」を追加「カレンダーイベントを取得」を追加「リマインダーを表示」を追加「アプリを開く」でNetEase Cloud Musicを選択。これだけです、すべてドラッグ＆ドロップで完了します。
\end{itemize}

実例2：自動「省電力モード」

\booksubsection{実例1：ワンタップ「おはようフロー」}

実例3：スキャン自動家計簿

\booksubsection{実例2：自動「省電力モード」}

もう一つ、知らない人が多い使い方：ショートカットをホーム画面に追加して、アプリのアイコンにできます。例えば「タクシーで帰宅」というショートカットを作り、アイコンをデスクトップに置けば、ワンタップでDiDiを開き、目的地を自宅に設定できます。アプリを開いて手動入力するより遥かに速いです。

\booksubsection{実例3：スキャン自動家計簿}

Android側は少し複雑ですが、上限もより高いです。まず簡単なものから：

入門レベル：オートクリッカー

\booksection{三、Androidオートメーション——Taskerとその仲間たち}

中級レベル：Tasker——Androidオートメーションの王

\booksubsection{入門レベル：オートクリッカー}

プロファイル（Profile）

\booksubsection{中級レベル：Tasker——Androidオートメーションの王}

タスク（Task）

\begin{itemize}
  \item \textbf{：「トリガー後に何をするか」を定義。一つのタスクは複数のアクションを含み、順番に実行されます。アクションの種類には、システム設定の調整（音量、明るさ、Wi-Fi、Bluetooth、機内モード）、SMS送信、HTTPリクエスト、Shellコマンドの実行、メディア再生の制御、ダイアログの表示などがあります。}Taskerの定番実用シナリオをいくつか紹介します：
  \item \textbf{帰宅時自動でWi-Fiオン}：トリガー条件を「特定の基地局に接続」（自宅付近では固定の基地局に接続するため）タスク「Wi-Fiをオンにする」。GPS不要、省電力で、スマートフォンが小区の範囲に入ると自動的にWi-Fiがオンになります。
\end{itemize}

会議モード

\begin{itemize}
  \item \textbf{：トリガー条件を「カレンダーイベントのタイトルに『会議』を含む」タスク「サイレントモード設定、メディア音量オフ、WeChat自動返信（プラグインが必要）」。会議中にスマートフォンが突然鳴る心配がなくなります。}夜間の目のケア
  \item \textbf{：トリガー条件を「時間22:00-06:00」タスク「アイケアモードをオン、明るさを20\%に下げる、画面を白黒（グレースケールモード）に設定」。白黒画面はショート動画を見る欲求に絶大な効果があります——画面がグレーになると、見る気が失せます。}イヤホン接続時自動ポッドキャスト
  \item \textbf{：トリガー条件を「有線イヤホン挿入/Bluetoothイヤホン接続」タスク「小宇宙アプリを開き、最新エピソードの再生を開始」。イヤホンを装着するだけで聴取でき、スマートフォンのロック解除してアプリを探す手間が省けます。}Taskerのハードルが高すぎると感じる場合は、簡易版の兄弟アプリ
  \item \textbf{MacroDroid}を試してみてください。インターフェースがはるか使いやすく、テンプレートも豊富で、こだわりたくないユーザーに適しています。また、現在多くの中国系Androidシステム（XiaomiのMIUI、HuaweiのHarmonyOS、OPPOのColorOS）に「シーン」や「スマートアシスタント」機能が標準搭載されており、Taskerほど強力ではありませんが、日常的なオートメーションには十分です。
\end{itemize}

四、オートメーション的思考の日常応用\textbf{オートメーションを学ぶ最大の収穫は、少しの時間を節約することではなく、}問題に対する考え方を変えること

\booksection{四、オートメーション的思考の日常応用}

写真を撮るたびに、自動的にクラウドストレージにバックアップし、ローカルの写真を小サイズに圧縮して容量を節約。\textbf{毎朝自動的にカレンダーから今日の予定を読み取り、テキストを生成してWeChatで家族に送信。}荷物の受領SMSを受信したら、自動的にメモに「荷物を受け取る」のToDoを追加。

\begin{itemize}
  \item バッテリーが満充電になったら自動的に音声で充電器を抜くようリマインド（過充電によるバッテリー劣化を防止）。
  \item 毎日の退勤時間になっても、まだ会社の近くにいる場合（基地局位置情報で判定）、自動的に「退勤の時間です」とリマインド。
  \item これらは高度に聞こえますが、「トリガー＋アクション」という2つの基本概念をマスターし、ショートカットやTaskerと組み合わせれば、すべて実現できます。重要なのは
  \item 「この作業は自動化できる」とまず気づくこと
  \item ——多くの人は「できない」のではなく、そういう発想すらないのです。
\end{itemize}

五、落とし穴ガイド：オートメーションのよくある問題\textbf{オートメーションも順風満帆ではありません。私が先に踏んだ落とし穴をいくつか紹介します：}位置トリガーが反応しない場合がある

\booksection{五、落とし穴ガイド：オートメーションのよくある問題}

ショートカットの「オートメーション」には認証が必要

\begin{itemize}
  \item \textbf{：iOSでは、「オートメーション」タブのトリガーの多くは、実行前に手動確認が必要です（iOSのセキュリティメカニズム）。ただし、Wi-Fi接続、Bluetooth接続、アラーム停止のトリガーは、確認なしで完全に自動実行できます。}Taskerはユーザー補助権限が必要
  \item \textbf{：多くの高度な機能（例えばタップのシミュレーション）にはユーザー補助サービス権限が必要で、設定画面に常駐通知が表示されます。これはAndroidのセキュリティ設計であり、削除できません。}オートメーションを複雑にしすぎない
  \item \textbf{：オートメーションを始める際は、「2-3のアクション」の小さなフローから始めてください。いきなり20ステップのスーパーフローを作ろうとしないでください——デバッグで発狂します。}ヒント：オートメーション入門の正しい姿勢
  \item \textbf{まず5分間考えてみてください：毎日スマートフォンで最も多く繰り返している3つの作業は何ですか？例えば毎朝のルーティン、Bluetoothイヤホン接続後の習慣、外出前の設定など……これらの高頻度の小さな作業から始めて、それぞれ簡単なオートメーションを作ってみてください。1つのオートメーションが5秒節約するだけでも、積み重なれば大きな節約になります。しかも、完成した時の達成感は、本当に癖になります。}練習：3つの入門オートメーションチャレンジ
\end{itemize}

\textbf{自分に最も適したものを1つ完了すれば、入門完了です：}\par

チャレンジ1：おはようオートメーション

\begin{tcbpractice}{：iOS：ショートカットアプリで新規ショートカット「現在の天気を取得」「テキストを朗読」「アプリを開くでNetEase Cloud Musicを選択」を追加。完了後、Siriに「おはよう」と言ってトリガーできます。Android：Taskerで新規タスク「天気を取得」「朗読」「アプリを開く」を追加トリガー時間を毎日7:00に設定。}
チャレンジ2：低バッテリー自動省電力

\begin{enumerate}
  \item \textbf{：iOS：オートメーションタブ新規「バッテリー残量」25\%を下回った時「低電力モードをオン」＋「明るさ25\%」を追加。Android：Taskerで新規Profile状態電源バッテリー0-25\%タスク「省電力モードオン」＋「明るさを下げる」。}チャレンジ3：カーブルートゥース接続時自動ナビ
  \item \textbf{：iOS：オートメーション新規「Bluetooth」カーブルートゥースを選択「アプリを開くでAmapを選択」＋「メディア音量を80\%に設定」を追加。Android：TaskerProfile状態Bluetooth接続済みカーブルートゥースを選択タスク「Amapを開く」＋「メディア音量を上げる」。}第二十三章付録：GitHub入門ガイド——コードホスティングからコラボレーションプラットフォームまで
  \item \textbf{前章で「プログラミング的思考」を説明しましたが、プログラミング的思考を実践し、他の人と一緒にプロジェクトを行うには、GitHubという名前を避けて通れません。コンピュータ専攻の学生がコンテストに参加する、自分のコード作品を保存したい、あるいは単に世界中の開発者が何に取り組んでいるか見たい場合でも、GitHubは必ず知っておくべきプラットフォームです。この節では最も基礎的な概念から説明し、明確な認識を構築します。}一、GitHubとは何か？
\end{enumerate}

\end{tcbpractice}

\begin{tcbinfo}{事例3：自動化「おはようフロー」}
\begin{itemize}
  \item \textbf{シーン}: 毎日7:00に自動で天気予報 + マナーモード解除 + ニュースを開く。
  \item \textbf{手順}: iOS: ショートカットオートメーション時間トリガー「天気を取得」「音量を設定」「Appを開く」を追加。Android: MacroDroidトリガー(時間 7:00)アクション(メディア音量/URLを開く/TTS天気読み上げ)。
  \item \textbf{期待される結果}: インストール不要で実現可能、導入ハードルが低い。
\end{itemize}

\end{tcbinfo}
